\section{从检测到预防：Intervention Ladder 与用户能力}

前面的系统大多在风险出现后识别和处置，本节把目标扩展为 prevention。预防不是“预测所有坏事”，而是在不同确定性下设计适当 friction：用户教育、默认隐私、年龄适配、举报入口、风险提示、联系限制、内容 hold 与专业升级。越不可逆的动作，所需证据和复核越高。

\subsection{四层 intervention ladder}

一个 practical ladder 可以从 educate、disrupt、protect 到 investigate/report 逐步升级。Education 帮助用户识别操纵和求助渠道；disruption 增加批量接触、陌生人消息或快速转移平台的成本；protection 提供默认设置、block/report 和 trusted-contact support；高置信事件才进入专业审核、平台处置与法定流程。

\lecturefigure{14-prevention-ladder.png}{预防需要教育、产品 friction、保护性升级和法定响应的分层组合。}{本地字幕 33:10--35:30；NCMEC NetSmartz 与课堂讨论；概念重绘。}

\begin{knowledgebox}{读图：Evidence 越弱，动作越应可逆}
低置信信号适合展示安全提示、限制陌生人发现或提供 block/report；中等风险可增加 friction、延迟发送或请求额外确认；高置信且经复核的事件再进入强制处置和报告。这个顺序保护用户免受误报，也让系统在早期仍能提供帮助。软件可以降低机会与响应时间，却不能替代家庭、学校、热线和专业服务。
\end{knowledgebox}

\teachervoice{Cordua 把方向概括为从 reactive detection 走向 prevention，但她没有宣称技术可以独自解决问题。家长和年轻人的数字素养、清晰求助渠道与平台默认保护仍是系统的一部分，而不是“模型上线后可以删除的人工环节”。}

\subsection{本章小结}

预防体系应把证据强度和动作可逆性匹配起来。Education、friction、user control、specialist review 与 lawful response 共同构成防线；任何单一 classifier 都不能覆盖这条 ladder。

